\documentclass[10pt,letterpaper]{article}

\usepackage[margin=0.45in]{geometry}
\usepackage[T1]{fontenc}
\usepackage{lmodern}
\usepackage{xcolor}
\usepackage{enumitem}
\usepackage{tabularx}
\usepackage[hidelinks]{hyperref}
\usepackage{titlesec}
\usepackage{microtype}

\definecolor{accent}{HTML}{176B70}
\setlength{\parindent}{0pt}
\setlength{\parskip}{0pt}
\setlist[itemize]{leftmargin=1.15em,itemsep=0.5pt,topsep=1pt,parsep=0pt}
\titleformat{\section}{\large\bfseries\color{accent}}{}{0pt}{}[\vspace{-2pt}\titlerule]
\titlespacing*{\section}{0pt}{4pt}{2pt}
\pagestyle{empty}

\newcommand{\role}[4]{%
  \textbf{#1}\hfill #2\\
  \textit{#3}\hfill \textit{#4}\vspace{-3pt}
}

\begin{document}
\fontsize{9}{10.35}\selectfont

\begin{center}
  {\Large\bfseries Piter Z. Garcia Bautista}\\[1pt]
  {\large Senior Data Analyst | SQL, Python, and R | Education and Research Workflows}\\[2pt]
  Rochester, NY \enspace|\enspace Remote\\
  \href{mailto:garciapiterz@gmail.com}{garciapiterz@gmail.com}
  \enspace|\enspace +1 (646) 288-3300
  \enspace|\enspace \href{https://linkedin.com/in/garciapiterz}{linkedin.com/in/garciapiterz}\\
  \href{https://github.com/pzg8794}{github.com/pzg8794}
  \enspace|\enspace \href{https://garciapiterz.com}{garciapiterz.com}
\end{center}

\section*{Summary}
Data analyst completing M.S. degrees in Data Science and Teaching Computer
Science K--12 in December 2026, with industry and research experience building
reliable data workflows. Experienced in SQL, Python, R, large and heterogeneous
datasets, data cleaning and validation, reproducible analysis, documentation,
and communicating findings to technical and nontechnical collaborators. Brings
K--12 teaching experience and a practical understanding of how data quality,
definitions, and limitations affect educators and learners.

\section*{Technical Skills}
\begin{tabularx}{\textwidth}{@{}>{\bfseries}l X@{}}
Programming & SQL, Python, R, Java, C++, C\#, JavaScript, Bash, MATLAB\\
Data workflows & pandas, NumPy, SciPy, data cleaning, transformation, validation, feature engineering, statistical analysis\\
Quality and reporting & Automated checks, structured logging, reproducible pipelines, technical documentation, Matplotlib, Plotly, Seaborn\\
Research and education & Research specifications, experiment evaluation, formative assessment, UDL-informed instruction, stakeholder communication\\
Systems & PostgreSQL, MySQL, AWS, GCP, Docker, Git/GitHub, Linux/Unix\\
\end{tabularx}

\section*{Relevant Experience}
\role{Graduate Assistant, Quantum and AI Research}{Aug 2025--Present}
{Rochester Institute of Technology, Software Engineering}{Rochester, NY}
\begin{itemize}
  \item Build Python analysis and evaluation workflows for adaptive decision
    systems using automated validation, structured logging, shared datasets,
    and reproducible experiment tracking.
  \item Translate research requirements and large experimental outputs into
    comparison-ready datasets, visualizations, technical reports, and
    manuscript-ready evidence.
\end{itemize}

\role{Data Solutions Engineer}{Sep 2022--Aug 2024}
{VEDADATA Inc.}{New York, NY}
\begin{itemize}
  \item Designed maintainable SQL, Python, and pandas workflows for data
    ingestion, cleaning, preprocessing, validation, and analytics-ready delivery.
  \item Built quality checks and diagnostics for missing, inconsistent, and
    unusual information; documented findings and data processes for technical
    and nontechnical collaborators.
\end{itemize}

\role{AI Data Solutions Engineer}{Jan 2021--Sep 2022}
{VIOME Inc.}{New York, NY}
\begin{itemize}
  \item Developed healthcare-oriented machine-learning workflows spanning
    preprocessing, feature engineering, model experimentation, evaluation,
    and deployment-oriented data preparation in AWS environments.
\end{itemize}

\role{Computer Science Teacher Candidate}{Feb 2026--May 2026}
{University of Rochester / Pine Brook Elementary School}{Rochester, NY}
\begin{itemize}
  \item Used formative checks and student participation evidence to adjust
    accessible K--5 computer-science instruction, grouping, pacing, and supports.
\end{itemize}

\role{Environmental Science Camp Teacher}{Jul 2026}
{University of Rochester / Freedom Scholars}{Rochester, NY}
\begin{itemize}
  \item Delivered inquiry-based science activities and adapted instruction in
    response to student understanding and engagement.
\end{itemize}

\section*{Selected Data and Research Work}
\textbf{Reproducible Analysis and Reporting.}
Built validation and reporting workflows that convert heterogeneous experiment
outputs into comparison-ready datasets, diagnostics, and reusable artifacts.

\textbf{Diabetic Readmission Prediction.}
Prepared and modeled a multi-hospital dataset using structured preprocessing,
feature engineering, model comparison, and performance evaluation.

\section*{Education}
\textbf{M.S., Data Science}, Rochester Institute of Technology
\hfill Expected December 2026\\
GPA: 3.9+; applied AI, bioinformatics, statistics, and reproducible data systems

\textbf{M.S., Teaching Computer Science K--12}, University of Rochester
\hfill Expected December 2026\\
GPA: 3.9+; NSF Robert Noyce Scholar; inclusive instruction and assessment

\textbf{M.S., Computer Science}, Rochester Institute of Technology
\hfill 2015\\
Artificial intelligence, big data analytics, and computer graphics

\textbf{Dual B.S., Computer and Electrical Engineering Technology}, Farmingdale State College
\hfill 2012

\end{document}
